\section{Penyangkalan dengan Kontra-Contoh}

Sejauh ini, pembahasan difokuskan pada cara \emph{membuktikan} bahwa suatu pernyataan benar. Namun, dalam matematika dan ilmu komputer, kemampuan untuk menunjukkan bahwa suatu pernyataan \emph{salah} sama pentingnya. Untuk menyangkal pernyataan universal (``untuk semua $x$, $P(x)$''), cukup menemukan satu elemen domain yang tidak memenuhi $P$ --- elemen tersebut disebut \logic{kontra-contoh} (\textit{counterexample}) \cite{rosen2019}.

\subsection{Landasan Logis}

Pernyataan universal $\forall x\, P(x)$ menyatakan bahwa $P(x)$ bernilai benar untuk setiap $x$ dalam domain. Negasi dari pernyataan ini adalah:
\[ \neg [\forall x\, P(x)] \equiv \exists x\, \neg P(x) \]

Artinya, untuk menyangkal pernyataan universal, cukup tunjukkan \emph{keberadaan} satu elemen $x_0$ dalam domain sehingga $P(x_0)$ bernilai salah. Elemen $x_0$ inilah yang disebut kontra-contoh.

\begin{definisi}[Kontra-Contoh]
\logic{Kontra-contoh} (\textit{counterexample}) adalah elemen spesifik dari domain yang menunjukkan bahwa suatu pernyataan universal adalah salah. Satu kontra-contoh sudah cukup untuk menyangkal pernyataan universal, tidak peduli berapa banyak contoh lain yang mendukungnya.
\end{definisi}

\subsection{Contoh Penyangkalan}

\begin{contoh}
\textbf{Klaim}: ``Untuk setiap bilangan bulat positif $n$, $n^2 - n + 41$ adalah bilangan prima.''

Klaim ini tampak benar jika diuji untuk nilai-nilai kecil:
\begin{center}
\renewcommand{\arraystretch}{1.2}
\begin{tabular}{|c|c|c|}
\hline
$n$ & $n^2 - n + 41$ & Prima? \\ \hline
1 & 41 & Ya \\ \hline
2 & 43 & Ya \\ \hline
3 & 47 & Ya \\ \hline
5 & 61 & Ya \\ \hline
10 & 131 & Ya \\ \hline
\end{tabular}
\end{center}

Namun, untuk $n = 41$:
\[ 41^2 - 41 + 41 = 41^2 = 1681 = 41 \times 41 \]
yang \textbf{bukan} bilangan prima. Kontra-contoh ditemukan pada $n = 41$, sehingga klaim tersebut \textbf{salah}. $\blacksquare$
\end{contoh}

\begin{contoh}
\textbf{Klaim}: ``Untuk setiap bilangan real $x$, $x^2 > x$.''

\textbf{Kontra-contoh}: Pilih $x = \frac{1}{2}$.
\[ x^2 = \left(\frac{1}{2}\right)^2 = \frac{1}{4} < \frac{1}{2} = x \]

Karena $x^2 < x$ (bukan $x^2 > x$), maka klaim tersebut \textbf{salah}. $\blacksquare$
\end{contoh}

\subsection{Peringatan: Kontra-Contoh vs Bukti}

\begin{konsep}
\textbf{Asimetri fundamental}: Satu kontra-contoh cukup untuk menyangkal pernyataan universal, tetapi seribu contoh yang mendukung \emph{tidak cukup} untuk membuktikan pernyataan universal. Membuktikan $\forall x\, P(x)$ memerlukan argumentasi yang berlaku untuk seluruh domain, bukan hanya sampel.

Analogi: Untuk membuktikan bahwa ``semua angsa berwarna putih'', memeriksa jutaan angsa putih tidak cukup. Tetapi menemukan satu angsa hitam sudah cukup untuk menyangkalnya.
\end{konsep}

Kemampuan membedakan antara verifikasi empiris dan bukti formal sangat penting dalam ilmu komputer. Pengujian perangkat lunak (\textit{testing}) menunjukkan keberadaan bug (mirip kontra-contoh), tetapi test yang lulus seluruhnya tidak membuktikan ketiadaan bug --- sama seperti contoh-contoh yang mendukung tidak membuktikan pernyataan universal \cite{velleman2019}.
